\begin{definitionbox}{Definition (Identity Relation)}
\begin{definition}[Identity Relation]
\label{def:identity-relation}
Let $A$ be a set. The \emph{identity relation} on $A$ is the binary relation
\[
\Delta_A = \{(a,a) \mid a \in A\}.
\]
\end{definition}

\begin{remark*}[Standard quantified statement]
\[
\operatorname{IdentityRelation}(\Delta_A)
\Longleftrightarrow
\forall x \in A\, \forall y \in A,\,
\bigl((x,y) \in \Delta_A \Longleftrightarrow x = y\bigr).
\]
\end{remark*}

\begin{remark*}[Interpretation]
The identity relation contains exactly the pairs in which an object is paired
with itself. It is the set-theoretic form of sameness inside a fixed domain.
\end{remark*}

\NoLocalDependencies
\end{definitionbox}

\begin{definitionbox}{Definition (Equality Relation)}
\begin{definition}[Equality Relation]
\label{def:equality-relation}
An \emph{equality relation} on a domain $A$ is the logical relation, written
$=$, that identifies precisely when two terms denote the same element of $A$.
\end{definition}

\begin{remark*}[Standard quantified statement]
\[
\operatorname{EqualityRelation}_A(=)
\Longleftrightarrow
\forall x \in A\, \forall y \in A,\,
(x = y \Longleftrightarrow x \text{ and } y \text{ denote the same element}).
\]
\end{remark*}

\begin{remark*}[Interpretation]
Equality is not an arbitrary similarity relation. It is the strict logical
relation of being the same object. Later constructions may introduce weaker
relations that group different objects together, but equality itself remains
the exact identity relation of the underlying domain.
\end{remark*}

\NoLocalDependencies
\end{definitionbox}

\begin{axiombox}{Axiom (Reflexivity of Equality)}
\begin{axiom}[Reflexivity of Equality]
\label{ax:equality-reflexivity}
Every object is equal to itself.
\[
\forall x \in A,\ x=x.
\]
\end{axiom}
\end{axiombox}

\begin{remark*}[Standard quantified statement]
\[
\operatorname{Reflexive}_A(=)
\Longleftrightarrow
\forall x \in A,\ x=x.
\]
\end{remark*}

\begin{remark*}[Interpretation]
Reflexivity is the irreducible identity principle: no object can fail to be
itself. This is the first logical constraint imposed on equality before any
arithmetic or order structure is introduced.
\end{remark*}

\NoLocalDependencies

\begin{axiombox}{Axiom (Substitution of Equals)}
\begin{axiom}[Substitution of Equals]
\label{ax:equality-substitution}
If two objects are equal, then replacing one by the other preserves every
well-formed property expressible in the language.
\[
\forall x,y \in A,\,
\Bigl(x=y \Longrightarrow
\bigl(\varphi(x) \Longleftrightarrow \varphi(y)\bigr)\Bigr).
\]
\end{axiom}
\end{axiombox}

\begin{remark*}[Standard quantified statement]
For every formula $\varphi(z)$ whose free occurrences of $z$ may be replaced
without variable capture,
\[
\forall x,y \in A,\,
\Bigl(x=y \Longrightarrow
\bigl(\varphi(x) \Longleftrightarrow \varphi(y)\bigr)\Bigr).
\]
\end{remark*}

\begin{remark*}[Interpretation]
Substitution is the operational content of equality. Equal objects are not
merely paired by a relation; they are interchangeable in every legitimate
mathematical assertion. Symmetry and transitivity are consequences of this
substitution behavior together with reflexivity.
\end{remark*}

\NoLocalDependencies

\begin{definitionbox}{Definition (Equivalence Relation)}
\begin{definition}[Equivalence Relation]
\label{def:equivalence-relation}
Let $A$ be a set. A binary relation $\sim$ on $A$ is an
\emph{equivalence relation} if it is reflexive, symmetric, and transitive.
\end{definition}

\begin{remark*}[Standard quantified statement]
\[
\operatorname{EquivalenceRelation}_A(\sim)
\Longleftrightarrow
\Bigl(\forall x \in A,\ x \sim x\Bigr)
\land
\Bigl(\forall x,y \in A,\ x \sim y \Rightarrow y \sim x\Bigr)
\land
\Bigl(\forall x,y,z \in A,\ (x \sim y \land y \sim z) \Rightarrow x \sim z\Bigr).
\]
\end{remark*}

\begin{remark*}[Interpretation]
An equivalence relation formalizes when objects should be treated as the same
for a chosen purpose, even if they are not literally equal. This distinction
is essential for quotient constructions: fractions, congruence classes, and
many algebraic structures are built by replacing exact identity with a
controlled equivalence relation.
\end{remark*}

\NoLocalDependencies
\end{definitionbox}
